\documentclass[10pt,a4paper]{amsart}
\usepackage[margin=19mm]{geometry}
\usepackage{amsmath,amssymb,mathtools}
\usepackage[scr=rsfs,cal=euler]{mathalfa}
\usepackage{xfrac}
\usepackage[unicode,hidelinks,bookmarks=false]{hyperref}
\newcommand{\ie}{i.e.\ }
\newcommand{\cf}{cf.\ }
\newcommand{\resp}{respectively\ }
\newcommand{\Subsection}{\subsection}
\providecommand{\nobreakdash}{\nobreak\hbox{-}}
\setlength{\emergencystretch}{2em}
\multlinegap0pt
\usepackage{mathtools}
\usepackage{amssymb}
\usepackage[shortlabels]{enumitem}
\usepackage{bm}
\newtheorem{proposition}{Proposition}
\newtheorem{lemma}{Lemma}
\usepackage{cleveref}
\crefname{proposition}{Proposition}{Propositions}
\crefname{lemma}{Lemma}{Lemmas}
\newcommand{\Prob}{\operatorname{Prob}}
\newcommand{\abs}[1]{\left\lvert#1\right\rvert}
\newcommand{\cN}{\mathcal N}
\newcommand{\cO}{\mathcal O}
\newcommand{\dbar}{\overline{\vol}}
\newcommand{\dd}{\,d}
\newcommand{\norm}[1]{\left\lVert#1\right\rVert}
\newcommand{\vol}{\operatorname{vol}}
\title{Poisson actions of noncompact locally compact sofic groups have completely positive entropy}
\author{Zhuowei Liu}
\date{Source: arXiv:2608.29599v1 (CC BY 4.0)}
\begin{document}
\maketitle

\noindent\emph{Source and licence.} Poisson actions of noncompact locally compact sofic groups have completely positive entropy. Version arXiv:2608.29599v1 (CC BY 4.0), distributed by arXiv under the Creative Commons Attribution 4.0 International licence, http://creativecommons.org/licenses/by/4.0/, as recorded in the arXiv licence metadata for this exact version and shown on the abstract page https://arxiv.org/abs/2608.29599v1. Full licence notice: \texttt{licenses/arxiv-2608.29599-CC-BY-4.0.txt}.\par\smallskip

\noindent\textit{Editorial scope.} Counted unit: the article's proposition prop:empirical-factor (source0.tex lines 1389--1394). The article's complete original proof of it is delivered with it (source0.tex lines 1396--1481, one \texttt{proof} environment of 86 lines). No mathematics is written, repaired or re-derived: the statement and the proof are the article's own lines, in the article's own notation, and no retained line is substituted or re-flowed.  Retained uncounted as the source-local context the proof needs: lines 1238--1254; lines 1346--1352.  Nothing was omitted from any retained span, so the package carries no omission record.  No retained line contains a citation, so the wrapper carries no bibliography and no bibliography entry.  No retained line contains a figure environment, an image-inclusion command or drawing code, so no image is delivered and the package declares no asset dependency.  Editorial formatting only, in addition: an AMS \texttt{amsart} preamble that restates the article's own declarations and macros used by the retained lines, namely the environments \texttt{proposition}, \texttt{lemma} on separate counters, together with the standard packages those lines need.  The retained environments print in this document as Proposition 1, Lemma 1 in order of appearance, whereas the article prints the same items with the numbering of its own full document.  The complete native source of the article as distributed in the arXiv e-print for this exact version is bundled unchanged as \texttt{sources/208868/source0.tex}.
\begin{proposition}[Concentration of local root averages]\label{prop:local-concentration}
Let $K\subset G$ be compact and let $F:\cN(G)\to[-1,1]$ be $K$-local. Define
\[
 \mathcal A_i(F):=\frac1{V_i}\int_{D_i}F(\Phi_i(p))\dd\vol_{M_i}(p).
\]
Then
\begin{equation}\label{eq:local-mean}
 \mathbb E\mathcal A_i(F)=d_i\int F\dd\mu_\alpha.
\end{equation}
Also, for every $t>0$ there is $c=c(t,K,\alpha)>0$ such that, for all sufficiently large $i$,
\begin{equation}\label{eq:local-concentration}
 \mathbb P\left(\abs{\mathcal A_i(F)-\mathbb E\mathcal A_i(F)}>t\right)
 \leq2e^{-cV_i}.
\end{equation}
The same claim holds for a function taking values in $[-1,1]$ that is
$K$-local in each coordinate of two independent completed root-view models.
\end{proposition}

\begin{lemma}\label{lem:local-L1}
For every bounded Borel function $H:\cN(G)\to\mathbb R$ and every $\eta>0$, there are a compact $K\subset G$ and a bounded $K$-local Borel function $H_K$ such that
\[
 \norm{H-H_K}_{L^1(\mu_\alpha)}<\eta,
 \qquad \norm{H_K}_\infty\leq\norm{H}_\infty.
\]
\end{lemma}

\begin{proposition}\label{prop:empirical-factor}
For every weak-star neighborhood $\cO$ of $\nu$, when $i \to \infty$,
\[
 \mathbb P\bigl((\Psi_i)_*\dbar_i\in\cO\bigr)\longrightarrow1.
\]
\end{proposition}

\begin{proof}
It is enough to treat a basic weak-star neighborhood
\[
 \cO=\left\{\lambda\in\Prob(Y):
 \left|\int f_j\,d\lambda-\int f_j\,d\nu\right|<\tau
 \text{ for }1\leq j\leq m\right\},
\]
where $f_j\in C(Y)$ and $\|f_j\|_\infty\leq1$. Put $H_j:=f_j\circ\pi$. To prove convergence in probability, fix an
arbitrary $\xi>0$. Choose $\eta>0$ so small that
\begin{equation}\label{eq:empirical-eta-choice}
 \sqrt\eta+2\eta<\tau/3,
 \qquad m\sqrt\eta<\xi/2.
\end{equation}
By \cref{lem:local-L1}, for each $j$ there is a compact window and a local approximation with $L^1$ error less than $\eta$. Replacing the finitely many windows by their union, we get one compact $K\subset G$ and $K$-local Borel functions $H_{j,K}$ satisfying
\[
 \|H_j-H_{j,K}\|_{L^1(\mu_\alpha)}<\eta,
 \qquad \|H_{j,K}\|_\infty\leq1
 \quad(1\leq j\leq m).
\]

Define
\[
 E_{i,j}:=\frac1{V_i}\int_{M_i}
 |H_j(\Phi_i(p))-H_{j,K}(\Phi_i(p))|\,d\vol_{M_i}(p).
\]
The exact marginal identity and Tonelli's theorem give
\[
 \mathbb E E_{i,j}
 =\int_{\cN(G)}|H_j-H_{j,K}|\,d\mu_\alpha<\eta.
\]
So we have
\begin{equation}\label{eq:empirical-approx-prob}
 \mathbb P(E_{i,j}>\sqrt\eta)<\sqrt\eta
\end{equation}
by Markov's inequality.

Set
\[
 J_{i,j}:=\frac1{V_i}\int_{D_i}H_{j,K}(\Phi_i(p))\,d\vol_{M_i}(p),
 \qquad m_{j,K}:=\int H_{j,K}\,d\mu_\alpha.
\]
By \cref{prop:local-concentration},
\[
 J_{i,j}-d_i m_{j,K}\xrightarrow{\mathbb P}0.
\]
Since $d_i\to1$, we get $J_{i,j}\to m_{j,K}$ in probability.
So there is $i_1$ such that, for every $i\geq i_1$,
\begin{equation}\label{eq:empirical-local-prob}
 \sum_{j=1}^m
 \mathbb P\left(|J_{i,j}-m_{j,K}|>\tau/3\right)<\xi/2.
\end{equation}
The contribution from the complement of the deep core satisfies
\[
 \left|\frac1{V_i}\int_{M_i\setminus D_i}
 H_{j,K}(\Phi_i(p))\,d\vol_{M_i}(p)\right|
 \leq1-d_i.
\]
Also,
\[
 |m_{j,K}-\int H_j\,d\mu_\alpha|<\eta.
\]
On the event $E_{i,j}\leq\sqrt\eta$,
\begin{align*}
 &\left|\frac1{V_i}\int_{M_i}H_j(\Phi_i(p))\,d\vol_{M_i}(p)
 -\int H_j\,d\mu_\alpha\right|\\
 &\quad\leq E_{i,j}+|J_{i,j}-m_{j,K}|+(1-d_i)+\eta.
\end{align*}
Increase $i_1$ if necessary so that
$1-d_i<\tau/3-\sqrt\eta-\eta$ for every $i\geq i_1$, which is possible by
\eqref{eq:empirical-eta-choice}. For such $i$, failure of any one of the $m$
required moment inequalities is contained in the union of the events
\[
 \{E_{i,j}>\sqrt\eta\}
 \quad\text{and}\quad
 \{|J_{i,j}-m_{j,K}|>\tau/3\},
 \qquad 1\leq j\leq m.
\]
By \eqref{eq:empirical-approx-prob},
\eqref{eq:empirical-local-prob}, and the union bound,
\[
 \mathbb P\bigl((\Psi_i)_*\dbar_i\notin\cO\bigr)
 \leq m\sqrt\eta+\xi/2<\xi.
\]
Here we used $\int H_j\,d\mu_\alpha=\int f_j\,d\nu$. Since $\xi>0$ was
arbitrary, the required probability tends to one.
\end{proof}

\end{document}
